\begin{frame}{copy et copy-of}
    \begin{block}{\textcolor{blue}{\texttt{<xsl:copy>}}}
    \begin{itemize}
        \item Copie que le noeud courant\\
        \textcolor{blue}{\texttt{<xsl:template \textcolor{orange}{match="body"}>}}\\
        \textcolor{blue}{<xsl:copy/>}\\
        \textcolor{blue}{</xsl:template>}
        
    \end{itemize}
    \end{block}

    \begin{block}{\textcolor{blue}{\texttt{<xsl:copy-of>}}}
        \begin{itemize}
            \item Copie le noeud courant et ses enfants:\\
            \textcolor{blue}{\texttt{<xsl:template \textcolor{orange}{match="body"}>}}\\
        \textcolor{blue}{<xsl:copy-of \textcolor{orange}{select="lg"}/>}\\
        \textcolor{blue}{</xsl:template>}
        \end{itemize}
    \end{block}
\end{frame}

\begin{frame}[fragile]{value-of}
    \begin{block}{\textcolor{blue}{\texttt{<xsl:value-of>}}}
    \begin{itemize}
        \item Permet d'extraire la valeur d'un élément.\\
        \begin{lstlisting}[numbers=none]
            <xsl:template match="lg">
                <xsl:value-of select="."/>
            </xsl:template>
        \end{lstlisting}    
    \end{itemize}
        
    \end{block}
\end{frame}

\begin{frame}[fragile]{number}
    \begin{block}{\textcolor{blue}{\texttt{<xsl:number>}}}
        \begin{itemize}
            \item Permet de produire un nombre dans l'arbre de résultat.
            \item Possède plusieurs attributs~:
            \begin{itemize}
                \item \textcolor{orange}{\texttt{@count}}: spécifier les noeuds à compter;
                \item \textcolor{orange}{\texttt{@level}}: spécifier le nombre de niveaux
                \begin{itemize}
                    \item \texttt{\textbf{any}}: compte le nombre de noeuds.
                    \item \texttt{\textbf{single}}: compte le nombre de noeuds descendants d'un même parent
                    \item \texttt{\textbf{multiple}}: compte tous les noeuds descendants mais aussi les "siblings" et leur descendant.
                \end{itemize}
                \item \textcolor{orange}{\texttt{@format}}: spécifier le format de numérotation.
            \end{itemize}
        \end{itemize}
        \begin{lstlisting}[numbers=none]
        <xsl:template match="mon_element_xml">
            <xsl:number count="element_a_compter" level="" format=""/>
        </xsl:template>
    \end{lstlisting}
    \end{block}
    
\end{frame}

\begin{frame}[fragile]{Exercice}
    \url{https://pad.libreon.fr/p9wWZK-ETAu6XIyAmOGjGg}
\end{frame}